\documentclass[10pt]{article}
\usepackage[utf8]{inputenc}
\usepackage[T1]{fontenc}
\usepackage{amsmath}
\usepackage{amsfonts}
\usepackage{amssymb}
\usepackage[version=4]{mhchem}
\usepackage{stmaryrd}

\begin{document}
где $w$ - произвольный $p$-мерный вектор, параллельный ребру многогранника $U$. Тогда для оптимальности по быстродействию управления $u(t), 0 \leqslant t \leqslant t_{1}$, переводящего объект из заданного начального состояния $x_{0}$ в положение равновесия (начало координат в пространстве $x$ ), необходимо и достаточно, чтобы оно удовлетворяло принципу максимума Понтрягина. Далеө, оптимальное управление $u(t)$ в линейной задаче оптимального быстродействия кусочно постоянно, и его значениями являются лишь вершины многогранника $U$.

В общем случае число переключений $u(t)$ хотя и ковечно, но может быть произвольным. В следующем важном случае число переключений допускает точную оценку сверху.

Если многогранник $U$ является $p$-мерным параллелепипедом

\[
a^{s} \leqslant u^{s} \leqslant b^{s}, \quad s=1, \ldots, p,
\]

и все собственные значения матрицы $\boldsymbol{A}$ действительны, то каждая из компонент $u^{s}(t), s=1, \ldots, p$, оптимального управления $u(t)$ является кусочно постоянной функцией, принимающей только значения $a^{s}$ и $b^{s}$ и имеющей не более $n-1$ переключениї, т. е. не более $n$ интервалов постоянства.

О. б. з. может рассматриваться и для неавтономных систем, т. е. для систем, у к-рых правая часть $f$ зависит еще и от времени $t$.

В тех случаях, когда әто удается, полезно рассматривать О.б. з. не только в программной постановке, как это описано выше, но и в позиционной постановке в форме задачи синтеза (см. Оптимальное управление позиционное). Решение задачі синтеза позволяет получить качественное представление о структуре оптимальвого по быстродействию управления, переводящего систему из любой точки, находящейся в нек-рой окрестности исходной начальной точки $x_{0}$, в заданное конечное состояние $x_{1}$.

Лum.: [1] П о н т р я ги н Л. С., Б о л т ян с к и й В. Г., Г а мк р е ли д з е Р. В., М и щ е н к о Е. Ф., Математическая с к и й В. Г., Математические методы оптимального управления, М., 1966.
И. Б. Вапнярский.

ОПТИМАЛЬНОГО УПРАВЛЕНИЯ МАТЕМАТИЧЕСКАЯ ТЕОРИЯ - раздел математики, в к-ром изучаются способы формализации и методы репения задач о выборе наилучшего в заранее предписанном смы̆сле способа осуществления управляемого динамич. процесса. Этот д и н а м и ч е с к й й п р о це с с может быть, как правило, описан при помощи дифференциальных, интегральных, функциональных, конечноразностных уравнений (или иных формализованных эволюционных соотношений), зависящих от системы функциїі пли параметров, наз. у п р а в л е н и я м и и подлежащих определению. Искомые управления, а также реализации самого процесса следует в общем случае выбирать с учетом ограничений, предписанных постановкой задачи.

В более специальном смысле термином «О. у. м. т.» принято называть математич. теорию, в к-рої изучаются методы решения неклассических вариационных задач оптимального управления (как правило, с дифференциальными связями), допускающих рассмотрение негладких функционалов и произвольных ограничений на параметры управления или иные зависимые переменные (обычно рассматривают ограничения, задаваемые нестрогими неравенствами). Термину «О. у. м. т.» иногда придают более широкий смысл, имея в виду теорию, изучающую математич. методы исследования задач, решения к-рых включают какой-либо процесс статической или динамич. оптимизации, а соответствующие модельные ситуации допускают интерпретацию в терминах той или иной прикладной процедуры принятия ваилучшего решения. В таком толковании О. у. м. т. содержит элементы исследования операций, математического программирования и игр теории.

Задачи, рассматриваемые в О. у. м. т., возникли из практич. потребностей, прежде всего в области механики космич. полета и автоматического управления теории (см. также Вариационное исчисление). Формализация п решение этих задач поставили новые вопросы, напр. в теории обыкновенных дифференциальных уравнений, как в области обобщения понятия решения и вывода соответствующих условий существования, так и в изучении динамических и әкстремальных свойств траекторий управляемых дифференциальных систем; в частности, О. у. м. т. стимулировала изучение свойств дифференциальных включений. Соответств ующиө направления О. у. м. т. поэтому часто рассматриваются как раздел төории обыкновенных диффференциальных уравнений. В О. у. м. т. содержатся математич. основы теории управления движением - нового раздела общей механики, в к-ром исследуются заковы формирования управляемых механич. движений и смежные математич. вопросы. По методам исследования и по своим приложениям О. у. м. т. тесно связана с аналитич. механикой, в особенности с разделами, относящимися к вариационным принципам классической механики.

Хотя частные задачи оптимального управления и неклассические вариационные задачи встречались и ранее, основы общей О. у. м. т. были заложены в $1956-$ 1961. Ключевым пунктом этой теории послужил Понтрягина принцип максим ума, сформулированный Л. С. Понтрягшным в 1956 (см. [1]). Важными стимулами создания О. у. м. т. были открытие метода динамического програмлирования, выяснение роли функционального анализа в теории оптимальных систем, открытие связей репений задач оптимального управления с результатами теориі устойчивости по Ляпунову, появление работ, связанных с понятиями управляемости и наблюдаемости динамич. систем (см. [2]-[5]). В последующие годы были развиты основы теории стохастич. управления іІ стохастич. фильтрации динамич. систем, построены общие методы решения неклассических вариационных задач, получены обобщения основных положений О. у. м. т. на более сложные классы динамич. систем, изучены связи с классическим вариационным исчислением (см. [6]-[11]). О. у. м. т. интенсивно развивается, в частности, в направлении изучения ігровых задач динамики (см. Дифференциальные игры), задач управления в условиях неполной или неопределенной информации, сшстем с распределенными параметрами, уравнений на многообразиях и т. Д.

Результаты О. у. м. т. нашли широкие приложения в формировании процессов управления, относящихся к самым разным областям сов ременной техники, в изучении экономич. динамики, в решении ряда задач из области бнологии, медицины, экологии, демографии и т. д.

3 а д ач а о птимального управ ления в общем виде может быть описана следующим образом.

\begin{enumerate}
  \item Дана управляемая систөма $S$, состояние к-рой в момент времени $t$ изображаөтся величиной $x$ (напр,, вектором обобщенных координат и обобщенных импульсов механич. системы, функцией от пространственных координат в распредөленной системе, вероятностным распределением, характеризующим текущее состояние стохастпч. системы, вектором выпуска продукции в динамич. модели әкономики и т. д.). Предполагается, что к системе $S$ приложены управляющие воздействия $u$, оказывающие влияние на ее динамику. Они могут, напр., иметь смысл механич. сил, температурных или электрич. потенциалов, программы капиталовложениї И т. Д.

  \item Дано уравнение, связывающее переменные $x$, $u, t$ II описывающие динамику системы. Указан про-

\end{enumerate}

\end{document}